\section{Entradas e Granularidade de Blocos}
\label{sec:entradas}

Para validar a corretude das implementações paralelas e avaliar diferentes regimes de casamento de caracteres e desvios condicionais, as entradas foram geradas sinteticamente via \texttt{generate\_input.py} combinando 5 padrões conceituais (Tabela~\ref{tab:padroes}).

\begin{table}[H]
\centering
\small
\caption{Padrões de caracteres e o score analítico esperado.}
\label{tab:padroes}
\begin{tabular}{@{}llll@{}}
\toprule
\textbf{Padrão} & \textbf{Sequência A} & \textbf{Sequência B} & \textbf{Score Esperado} \\ \midrule
\texttt{p1} & \texttt{"abcd"} repetido & \texttt{"bcda"} repetido & $\min(N,M) - 1$ \\
\texttt{p2} & \texttt{"a"} repetido & \texttt{"b"} repetido & $0$ \\
\texttt{p3} & \texttt{"abcd"} repetido & \texttt{"defg"} repetido & $\min(N,M)/4$ \\
\texttt{p4} & \texttt{"triviall"} repetido & \texttt{"triviall"} repetido & $\min(N,M)$ \\
\texttt{random} & Bytes pseudo-aleatórios & Bytes pseudo-aleatórios & $\approx 0{,}84 \min(N,M)$ (empírico) \\ \bottomrule
\end{tabular}
\end{table}

\subsection{Famílias de Matrizes e Parâmetros dos Ladrilhos}

A Tabela~\ref{tab:familias_entradas} sintetiza o conjunto de matrizes avaliadas, abrangendo diferentes volumes, proporções de aspecto e restrições de memória.

\begin{table}[H]
\centering
\footnotesize
\caption{Famílias de entradas, dimensões, padrões de caracteres e modos avaliados.}
\label{tab:familias_entradas}
\resizebox{\textwidth}{!}{%
\begin{tabular}{@{}llcc@{}}
\toprule
\textbf{Família / Casos} & \textbf{Dimensões ($M \times N$)} & \textbf{Padrões Avaliados} & \textbf{Modos Suportados} \\ \midrule
enunciado / quebrado & $10 \times 7$ / $827 \times 827$ & fixo / \texttt{p1} & \texttt{--full-matrix}, \texttt{--mem-opt} \\
Quadradas: Q1 a Q6 & $1.000^2$ até $100.000^2$ & \texttt{p1}, \texttt{p2}, \texttt{p3}, \texttt{p4}, \texttt{random} & \texttt{--full-matrix}, \texttt{--mem-opt} \\
Quadradas: Q7 e Q8 & $160.000^2$ e $320.000^2$ & Q7: todos (5); Q8: \texttt{p1} & apenas \texttt{--mem-opt} (giga) \\
Retangulares: R1 a R3 & $80k \times \{40, 20, 10\}k$ (e inv.) & todos (5) por variante & \texttt{--full-matrix}, \texttt{--mem-opt} \\
Retangulares: R4 a R6 & $160k \times \{80, 40, 20\}k$ (e inv.) & \texttt{p1} (direto e inv.) & apenas \texttt{--mem-opt} (giga) \\ \bottomrule
\end{tabular}%
}
\end{table}

Quanto à granularidade dos blocos, adotou-se \texttt{BLOCK\_SIZE} = \textbf{512} no OpenMP (definido empiricamente na Seção~\ref{sec:blocksize}) e \texttt{CHUNK\_SIZE} = \textbf{1024} no MPI (escolha de projeto para balancear latência e vazão no pipeline).
